\subsection{2  Лекция 2: Классификация уравнений}

Итак, чтобы упростить уравнение (ЛДУЧП2ППК) с помощью неособенной линейной замены независимых переменных ${\rm \xi }=S^{T} x$, $\det S\ne 0$, достаточно упростить квадратичную форму $K=\sum _{i,j=1}^{n}a_{ij} y_{i} y_{j}  =y^{T} Ay$ с помощью неособенной линейной замены переменных $y=S{\rm \eta }$. Из курса линейной алгебры известно, что всегда существует неособенная линейная замена переменных $y=S{\rm \eta }$, при которой квадратичная форма $K$ принимает следующий канонический вид:
\[K=\sum _{k=1}^{p}{\rm \eta }_{k}^{2}  -\sum _{k=p+1}^{p+q}{\rm \eta }_{k}^{2}  ={\rm \eta }_{1}^{2} +...+{\rm \eta }_{p}^{2} -{\rm \eta }_{p+1}^{2} -...-{\rm \eta }_{p+q}^{2} \]
где $p$ и $q$ -- соответственно положительный и отрицательный индексы инерции квадратичной формы $K$, при этом $r=p+q$ -- ранг квадратичной формы.

Тогда соответствующая неособенная линейная замена независимых переменных ${\rm \xi }=S^{T} x$ приводит уравнение (ЛДУЧП2ППК) к следующему каноническому виду:
\[\tilde{L}\tilde{u}=\sum _{k=1}^{p}\frac{\partial ^{2} \tilde{u}}{\partial {\rm \xi }_{k}^{{\rm 2}} }  -\sum _{k=p+1}^{p+q}\frac{\partial ^{2} \tilde{u}}{\partial {\rm \xi }_{k}^{{\rm 2}} }  +\sum _{k=1}^{n}\tilde{b}_{k} \frac{\partial \tilde{u}}{\partial {\rm \xi }_{k} }  +c\tilde{u}=\tilde{f}({\rm \xi })\]

В силу закона инерции квадратичных форм значения положительного $p$ и отрицательного $q$ индексов инерции квадратичной формы не зависят от неособенной линейной замены переменных $y=S{\rm \eta }$. Это позволяет классифицировать дифференциальные уравнения (ЛДУЧП2ППК) в зависимости от значений коэффициентов $a_{ij} $ ($i,j=1,...,n$) при старших (вторых) производных.

1. Если в квадратичной форме $K=y^{T} Ay$ ранг $r=n$ и в её каноническом виде все слагаемые одного знака (т.е. либо $p=r=n$ ($q=0$), либо $p=0$ ($q=r=n$)), то уравнение (ЛДУЧП2ППК) называется уравнением \textit{эллиптического} типа. Например, уравнение Лапласа -- уравнение эллиптического типа:
\[\frac{\partial ^{2} u}{\partial x^{2} } +\frac{\partial ^{2} u}{\partial y^{2} } +\frac{\partial ^{2} u}{\partial z^{2} } =0\]

2. Если в квадратичной форме $K=y^{T} Ay$ ранг $r=n$, но в её каноническом виде имеются слагаемые разных знаков (т.е. $1\le p\le n-1$ и $1\le q\le n-1$), то уравнение (ЛДУЧП2ППК) называется уравнением \textit{гиперболического} типа (при $p=1$ ($q=n-1$) или при $p=n-1$ ($q=1$) -- \textit{нормально гиперболического} типа). Например, волновое уравнение -- уравнение гиперболического типа:
\[\frac{\partial ^{2} u}{\partial t^{2} } =c^{2} \left(\frac{\partial ^{2} u}{\partial x^{2} } +\frac{\partial ^{2} u}{\partial y^{2} } +\frac{\partial ^{2} u}{\partial z^{2} } \right)\]

3. Если в квадратичной форме $K=y^{T} Ay$ ранг $r<n$, то уравнение (ЛДУЧП2ППК) называется уравнением \textit{параболического} типа (при $r=n-1$ и $p=0$ ($q=n-1$) или $p=n-1$ ($q=0$) -- \textit{нормально параболического} типа). Например, уравнение теплопроводности -- уравнение параболического типа:
\[\frac{\partial u}{\partial t} =a^{2} \left(\frac{\partial ^{2} u}{\partial x^{2} } +\frac{\partial ^{2} u}{\partial y^{2} } +\frac{\partial ^{2} u}{\partial z^{2} } \right)\]

Вернемся к рассмотрению в области $G\subset R^{n} $ \textbf{дифференциального уравнения с частными производными второго порядка, линейного относительно старших производных -- производных второго порядка (ДУЧП2ПЛОСП):}
\[\sum _{i,j=1}^{n}a_{ij} (x)\frac{\partial ^{2} u}{\partial x_{i} \partial x_{j} }  +\Phi \left(x,u,{\rm grad}u\right)=0\]
где коэффициенты при старших (вторых) производных $a_{ij} (x)\in C(G)$ ($i,j=1,...,n$) являются функциями только независимых переменных $x_{1} $,{\dots},$x_{n} $, и образуют квадратную матрицу $A(x)=(a_{ij} (x))$ -- \textit{матрицу старших коэффициентов}. Как было отмечено ранее, не умаляя общности, будем считать матрицу $A(x)$ симметрической ($A^{T} (x)=A(x)$, $a_{ij} (x)=a_{ji} (x)$).

\noindent

Выясним, по какому закону преобразуются коэффициенты $a_{ij} $ ($i,j=1,...,n$) при произвольной неособенной замене независимых переменных ${\rm \xi }={\rm \xi }(x)$, т.е.
\[{\rm \xi }_{k} ={\rm \xi }_{k} (x_{1} ,...,x_{n} )\in C^{2} (G),   k=1,...,n,   \frac{\partial {\rm (\xi }_{1} {\rm ,...,\xi }_{n} {\rm )}}{\partial (x_{1} ,...,x_{n} )} \ne {\rm 0}\]

Так как якобиан не равен нулю, то в некоторой окрестности можно выразить переменные $x_{1} ,...,x_{n} $ через переменные ${\rm \xi }_{1} {\rm ,...,\xi }_{n} $: $x=x({\rm \xi })$. Обозначим $u(x({\rm \xi }))=\tilde{u}({\rm \xi })$, тогда $\tilde{u}({\rm \xi (}x{\rm )})=u(x)$. Имеем

\noindent
\[\frac{\partial u}{\partial x_{i} } =\sum _{k=1}^{n}\frac{\partial \tilde{u}}{\partial {\rm \xi }_{k} } \frac{\partial {\rm \xi }_{k} }{\partial x_{i} }  ,   \frac{\partial u}{\partial x_{j} } =\sum _{l=1}^{n}\frac{\partial \tilde{u}}{\partial {\rm \xi }_{l} } \frac{\partial {\rm \xi }_{l} }{\partial x_{j} }  \]
\[\frac{\partial ^{2} u}{\partial x_{i} \partial x_{j} } =\frac{\partial }{\partial x_{i} } \left(\frac{\partial u}{\partial x_{j} } \right)=\sum _{k,l=1}^{n}\frac{\partial ^{2} \tilde{u}}{\partial {\rm \xi }_{k} \partial {\rm \xi }_{l} } \frac{\partial {\rm \xi }_{k} }{\partial x_{i} } \frac{\partial {\rm \xi }_{l} }{\partial x_{j} }  +\sum _{k=1}^{n}\frac{\partial \tilde{u}}{\partial {\rm \xi }_{k} } \frac{\partial ^{2} {\rm \xi }_{k} }{\partial x_{i} \partial x_{j} }  \]
или, в операторном виде
\[\frac{\partial }{\partial x_{i} } =\sum _{k=1}^{n}\frac{\partial {\rm \xi }_{k} }{\partial x_{i} } \frac{\partial }{\partial {\rm \xi }_{k} }  ,  \frac{\partial ^{2} }{\partial x_{i} \partial x_{j} } =\sum _{k,l=1}^{n}\frac{\partial {\rm \xi }_{k} }{\partial x_{i} } \frac{\partial {\rm \xi }_{l} }{\partial x_{j} } \frac{\partial ^{2} }{\partial {\rm \xi }_{k} \partial {\rm \xi }_{l} }  +\sum _{k=1}^{n}\frac{\partial ^{2} {\rm \xi }_{k} }{\partial x_{i} \partial x_{j} } \frac{\partial }{\partial {\rm \xi }_{k} }  \]
Подставляя найденные выражения в уравнение (ДУЧП2ПЛОСП), получим
\[\sum _{k,l=1}^{n}\frac{\partial ^{2} \tilde{u}}{\partial {\rm \xi }_{k} \partial {\rm \xi }_{l} } \sum _{i,j=1}^{n}a_{ij} \frac{\partial {\rm \xi }_{k} }{\partial x_{i} } \frac{\partial {\rm \xi }_{l} }{\partial x_{j} }   +\sum _{k=1}^{n}\frac{\partial \tilde{u}}{\partial {\rm \xi }_{k} } \sum _{i,j=1}^{n}a_{ij} \frac{\partial ^{2} {\rm \xi }_{k} }{\partial x_{i} \partial x_{j} }   +\hat{\Phi }\left({\rm \xi },\tilde{u},{\rm grad}\tilde{u}\right)=0\]

Обозначим теперь через где $\tilde{a}_{kl} $ ($k,l=1,...,n$) новые коэффициенты при вторых производных:
\[\tilde{a}_{kl} ({\rm \xi })=\sum _{i,j=1}^{n}a_{ij} (x({\rm \xi }))\frac{\partial {\rm \xi }_{k} }{\partial x_{i} } \frac{\partial {\rm \xi }_{l} }{\partial x_{j} }  \]
перепишем полученное уравнение в виде:
\[\sum _{k,l=1}^{n}\tilde{a}_{kl} ({\rm \xi })\frac{\partial ^{2} \tilde{u}}{\partial {\rm \xi }_{k} \partial {\rm \xi }_{l} }  +\tilde{\Phi }\left({\rm \xi },\tilde{u},{\rm grad}\tilde{u}\right)=0\]
где
\[\tilde{\Phi }\left({\rm \xi },\tilde{u},{\rm grad}\tilde{u}\right)=\sum _{k=1}^{n}\frac{\partial \tilde{u}}{\partial {\rm \xi }_{k} } \sum _{i,j=1}^{n}a_{ij} \frac{\partial ^{2} {\rm \xi }_{k} }{\partial x_{i} \partial x_{j} }   +\hat{\Phi }\left({\rm \xi },\tilde{u},{\rm grad}\tilde{u}\right)\]

Фиксируем точку $x^{(0)} =(x_{1}^{(0)} ,...,x_{1}^{(0)} )\in G$; обозначим ${\rm \xi }^{(0)} ={\rm \xi }(x^{(0)} )$,
\[s_{ik} =\frac{\partial {\rm \xi }_{k} (x^{(0)} )}{\partial x_{i} }    i,k=1,...,n.\]
Тогда формула для $\tilde{a}_{kl} ({\rm \xi })$ в точке $x^{(0)} $ запишется в виде:
\[\tilde{a}_{kl} ({\rm \xi }^{(0)} )=\sum _{i,j=1}^{n}a_{ij} (x^{(0)} )s_{ik} s_{jl}  ,\]
или, иначе (в матричном виде),
\[\tilde{A}({\rm \xi }^{(0)} )=S^{T} A(x^{(0)} )S\]
где $\tilde{A}({\rm \xi }^{(0)} )=(\tilde{a}_{kl} ({\rm \xi }^{(0)} ))$, $S=(s_{ik} )$.

Полученная формула преобразования коэффициентов $a_{ij} $ ($i,j=1,...,n$) в точке $x^{(0)} $ совпадает с формулой преобразования коэффициентов квадратичной формы

\noindent $\displaystyle K=\sum _{i,j=1}^{n}a_{ij} (x^{(0)} )y_{i} y_{j}  =y^{T} A(x^{(0)} )y$,   где    $\displaystyle y=\left(\begin{array}{c} {y_{1} } \\ {...} \\ {y_{n} } \end{array}\right)$

\noindent при неособенной линейной замене переменных

\noindent $\displaystyle y_{i} =\sum _{k=1}^{n}s_{ik} {\rm \eta }_{k}  $   (или, иначе, $y=\left(\begin{array}{c} {y_{1} } \\ {...} \\ {y_{n} } \end{array}\right)=S\left(\begin{array}{c} {{\rm \eta }_{1} } \\ {...} \\ {{\rm \eta }_{n} } \end{array}\right)=S{\rm \eta }$),   $\det S\ne 0$,

\noindent приводящей форму $K$ к виду:
\[K=\sum _{k,l=1}^{n}\tilde{a}_{kl} ({\rm \xi }^{(0)} ){\rm \eta }_{k} {\rm \eta }_{l}  ={\rm \eta }^{T} \tilde{A}({\rm \xi }^{(0)} ){\rm \eta }\]
